%!TEX root = ../template.tex
%%%%%%%%%%%%%%%%%%%%%%%%%%%%%%%%%%%%%%%%%%%%%%%%%%%%%%%%%%%%%%%%%%%%
%% chapter5.tex
%% NOVA thesis document file
%%
%% Chapter with lots of dummy text
%%%%%%%%%%%%%%%%%%%%%%%%%%%%%%%%%%%%%%%%%%%%%%%%%%%%%%%%%%%%%%%%%%%%

\typeout{NT FILE chapter5.tex}%

\chapter{Parametope}
\label{ch:parametope}

Several representations of reachable sets of dynamical have been introduced earlier, as well as methods of propagating them. However, they all have their pitfalls, especially with non-linear systems. The study ~\cite{harapanahalli2025parametric} offers a new framework for parametric reachable set computation through dynamical embeddings that demonstrates good accuracy, scalability and computational tractability. \\

It introduces a new set representation in the form of the parametope, represented as

\begin{equation*}
    ((\mathring{x}, \alpha, y))_g := \{ x \in \mathbb{R}^n : g(\alpha, x-\mathring{x}) \le y \}
\end{equation*}

where $\mathring{x} \in \mathbb{R}^n$ (center), $\alpha \in \mathbb{R}^p$ (parameters), $y \in \mathbb{R}^K$ (offset) and $g : \mathbb{R}^m \times \mathbb{R}^n \rightarrow \mathbb{R}^K$ (non-linearity). Different types of sets may be represented using parametopes, such as H-rep polytopes, ellipsoids, and annulus'. \\

The study further provides a dynamical framework for reachable
set computation by embedding the original system into a new
system evolving on $\mathbb{R}^{n} \times \mathbb{R}^{p} \times \mathbb{R}^{K}$, referred to as controlled embeddings. The embedding system is defined as

\begin{equation*}
    \dot{\mathring{x}} = f(\mathring{x}, \mathring{w}), \quad \dot{\alpha} = U(\mathring{x}, \alpha, y), \quad \dot{y} = E(\mathring{x}, \alpha, y),
\end{equation*}

evolving on $\mathbb{R}^{n} \times \mathbb{R}^{p} \times \mathbb{R}^{K}$, where $U$ is an arbitrary locally Lipschitz map, and $E$ is a locally Lipschitz map satisfying $ (x,w) \in ((\mathring{x}, \alpha, y))_g \times \mathcal{W} \Rightarrow E(\mathring{x}, \alpha, y) \le \xi(t,x,w)$. The term $\xi$ characterises the infinitesimal change in offset for any $x$, and is defined as
\begin{equation*}
    \xi (t,x,w) := \frac{\partial g}{\partial \alpha}(\alpha, x, \mathring{x} ) [ U(\mathring{x}, \alpha, y) ] + \frac{\partial g}{\partial x} (\alpha, x - \mathring{x}) [ f(x,w) - f(\mathring{x}, \mathring{w}) ]. \\
\end{equation*}

The study's Example 3 utilizes controlled ellipsoids to compute the reachable set of a robot arm. The example is replicated here, adapted to the MATLAB environment, achieving similar results. \\

Consider an autonomous system $\dot{x} = f(x)$. Furthermore, $g(\alpha , x - \mathring{x}) = ||\alpha(x-\mathring{x})||_2^2$, where $\alpha$ is a square invertible matrix. With the ellipsoid reachable set, the error dynamics of the system are overapproximated with a linear differential inclusion, this is, $f(x)-f(\mathring{x}) \in \mathcal{M} (x-\mathring{x})$, where $\mathcal{M} \subseteq \mathbb{R}^{n \times n}$, where $\mathcal{M}$ is a set of matrices (aka tensor?). In this case it is obtained by building the interval mixed Jacobian (IMJ) or mixed Jacobian LDI. For this end, the MATLAB function \verb|getElements| was built, which recursively breaks apart a symbolic expression into a binary tree of sub-expressions, where the leaf nodes are individual variables or constants, and calculates the range of the expression, leaf to stem, with interval arithmetic. The \verb|updateElement| function recalculates the range of an existing range tree when the bounds of the variables change. The IMJ is obtained by applying these functions to each non constant element of the system's Jacobian. \\

The embedding system is further defined by $U = -\alpha \text{Df}(\mathring{x})$ and $E = cy$, where $c = \underset{i}{\text{max}} \lambda_{max} ((\alpha (M_i - \text{Df}(\mathring{x})) \alpha^{-1})^T + \alpha (M_i - \text{Df}(\mathring{x})) \alpha^{-1})$, and $\mathcal{M} \subseteq \text{co}\{M_i\}_i$. The $c$ constant is calculated with the function \verb|prmtpEllips_c| which takes a brute force approach to find the maximum eigenvalues. This is, of course, impractical for larger systems. A possible optimization would be to take a gradient descent-based approach. \\

\begin{figure}[htbp]
  \centering
  \includegraphics[width=0.8\linewidth]{parametope_example3.png}%
  \caption{Results of parametope propagation}
  \label{fig:parametope_example3}
\end{figure}

